% Irodov VI §6.3 original 118--134.
% Printed pages 246--247, supplied PDF pp.255--256, answers pp.345--346.
\atomictermsproblem{118}{246}
某元素原子的K、L、M三主壳层及$4s$副壳层均全满，
而$4p$副壳层恰半满，求元素名称与原子序数。
\exsolution
封闭K、L、M壳层电子数分别为$2,8,18$；
$4s$全满再贡献2个电子；
$4p$可容6个，半满即3个。
总电子数为
\[
 Z=2+8+18+2+3=\boxed{33}.
\]
故该元素为$\boxed{\mathrm{As}\text{（砷）}}$，
构型简记为$[\mathrm{Ar}]3d^{10}4s^24p^3$。
\excheck 题中M壳全满含$3d^{10}$，
不可只数$3s^23p^6$。

\atomictermsproblem{119}{246}
用洪特规则求只有一个未满$p$副壳层、
分别含(a)3个和(b)4个电子的原子基态谱项。
\exsolution
$p$有$l=1$及三个$m_l=-1,0,1$。
(a) $p^3$为半满，三个电子自旋平行，
$S=3/2$且$L=0$，故$J=3/2$，
\[
 \boxed{p^3:\ {}^{4}S_{3/2}}.
\]
(b) $p^4$等价于两个$p$空穴，
高自旋$S=1$，合轨道$L=1$；
超过半满时$J=L+S=2$，
\[
 \boxed{p^4:\ {}^{3}P_2}.
\]
\excheck 洪特第三规则：少于半满取$J=|L-S|$，
超过半满取$J=L+S$。

\atomictermsproblem{120}{246}
用洪特规则求仅有未满$d$副壳层且其电子数
分别为(a)3、(b)7时原子的基态总角动量模。
\exsolution
$d$轨道$l=2$，
对$d^3$有三自旋同向$S=3/2$，
最大$L=2+1+0=3$（F）。
少于半满取$J=|3-3/2|=3/2$，
\[
 \boxed{d^3:\ {}^{4}F_{3/2},\quad
 M_J=\frac{\sqrt{15}}2\hbar}.
\]
对$d^7$可用三空穴等效，
$S=3/2,L=3$；超过半满取
$J=3+3/2=9/2$，
\[
 \boxed{d^7:\ {}^{4}F_{9/2},\quad
 M_J=\frac{\sqrt{99}}2\hbar}.
\]
\excheck 总角动量模是$\hbar\sqrt{J(J+1)}$，
不是$J\hbar$；$d^7$按多于半满处理。

\atomictermsproblem{121}{246}
一个原子只有一个未满副壳层，基态分别为
(a)\({}^{3}F_2\)；
(b)\({}^{2}P_{3/2}\)；
(c)\({}^{6}S_{5/2}\)。
用洪特规则求该副壳层电子数。
\exsolution
(a) $F$态$L=3$、三重自旋$S=1$，
$J=|L-S|=2$表明少于半满；
可由两个$d$电子构成，
故$\boxed{d^2}$。
(b) $P$态$L=1,S=1/2$，
$J=L+S=3/2$，相应$p$超过半满，
故$\boxed{p^5}$。
(c) $S$态$L=0,S=5/2$，
来自五个自旋同向$d$电子的半满壳层，
故$\boxed{d^5}$。
\excheck 上述指唯一未满副壳层，
不包含闭合壳层内电子数。

\atomictermsproblem{122}{246}
用洪特规则确定唯一未满副壳层原子的基态：
(a)壳层装满了最大容量的$1/3$，且$S=1$；
(b)装满$70\%$，且$S=3/2$。
\exsolution
(a) 对$p$副壳层最大6个位置，
装入2个即占$1/3$，可得
$S=1,L=1$；少于半满取$J=0$，
\[
 \boxed{{}^{3}P_0}.
\]
(b) 对$d$副壳层最大10个位置，
装入7个即占$70\%$，其
$S=3/2,L=3$，超过半满取$J=9/2$，
\[
 \boxed{{}^{4}F_{9/2}}.
\]
\excheck 含电子数必须为整数，
同时满足指定的自旋量子数和填充分数。

\atomictermsproblem{123}{246}
一个原子仅有一个未满副壳层，
该副壳层有三电子，
其基态的轨道量子数$L=3$。
用洪特规则写出基态谱项。
\exsolution
$p^3$的最大轨道量子数为$L=0$，
因此$L=3$表明三电子应在$d$副壳层。
三个$d$电子自旋同向$S=3/2$，
最大轨道$L=3$。
少于半满，$J=|L-S|=3/2$，
\[
 \boxed{{}^{4}F_{3/2}}.
\]
\excheck 若先判定该电子构型为$d^3$，
再按洪特第三规则即可确定下标$J$。

\atomictermsproblem{124}{247}
一个原子的唯一未满副壳层恰含五电子且半满。
求其基态磁矩的模（以玻尔磁子表示）。
\exsolution
半满五电子只能为$d^5$，
洪特规则给$S=5/2,L=0$，
因而$J=5/2$，基态谱项为\({}^{6}S_{5/2}\)。
纯S态朗德因子$g=2$，
\[
 \boxed{\mu=g\mu_B\sqrt{J(J+1)}
 =2\mu_B\sqrt{(5/2)(7/2)}
 =\sqrt{35}\,\mu_B\simeq5.92\mu_B}.
\]
\excheck 此磁矩是向量模，不是沿外场的最大投影
$gJ\mu_B=5\mu_B$。

\atomictermsproblem{125}{247}
温度$T=3000\,\mathrm K$时，
有多少比例的氢原子处于主量子数$n=2$态？
假定稀薄气体满足热平衡且几乎全部在基态。
\exsolution
氢能级$E_n=-13.6/n^2\,\mathrm{eV}$，
故$1\to2$激发能为$\Delta E=10.2\,\mathrm{eV}$。
忽略微小精细结构，
$n$级统计权重与$n^2$成正比，
$g_2/g_1=4$。
玻尔兹曼比例
\[
 \boxed{\frac{N_2}{N_1}=
 4e^{-10.2\,\mathrm{eV}/(k_BT)}
 \simeq2.9\times10^{-17}}.
\]
\excheck 本题假设中性氢并处于热平衡；
如此极小的激发份额可近似为总原子中的比例。

\atomictermsproblem{126}{247}
钠蒸气温度$T=2400\,\mathrm K$。
求处于$3P$与基态$3S$的原子数比。
已知$3P\to3S$共振线波长$\lambda=589\,\mathrm{nm}$。
\exsolution
$3P$的轨道角动量$L=1,S=1/2$，
包括6个磁-自旋态；$3S$仅2个，
故$g_P/g_S=6/2=3$。
激发能$\Delta E=hc/\lambda$，
玻尔兹曼分布给
\[
 \boxed{\frac{N_{3P}}{N_{3S}}=
 3\exp\!\left[-\frac{hc}{\lambda k_BT}\right]
 \simeq1.14\times10^{-4}}.
\]
\excheck 必须保留简并度因子3，
不能单用$e^{-\Delta E/kT}$。

\atomictermsproblem{127}{247}
激发原子形成的定向束流速率$v=600\,\mathrm{m/s}$。
沿束流前进$l=2.5\,\mathrm{mm}$后，
通往基态的谱线强度降为原来的$1/\eta$，
其中$\eta=25$。
求平均激发态寿命。
\exsolution
激发态数目满足
$N(t)=N(0)e^{-t/\tau}$；
自发辐射强度正比于$N$。
飞行时间$t=l/v$，因而
$e^{-l/(v\tau)}=1/\eta$，
\[
 \boxed{\tau=\frac l{v\ln\eta}
 =\frac{2.5\times10^{-3}}
 {600\ln25}\,\mathrm s\simeq1.3\,\mu\mathrm s}.
\]
\excheck 沿程衰减是激发态统计指数衰减，
不是电磁波在介质中的比尔衰减。

\atomictermsproblem{128}{247}
稀薄汞蒸气原子本来几乎全部处于基态，
经共振波长$\lambda=253.65\,\mathrm{nm}$的汞灯照射。
蒸气以该波长再辐射功率$P=35\,\mathrm{mW}$，
共振态平均寿命$\tau=0.15\,\mu\mathrm s$。
求共振激发态原子数。
\exsolution
每个激发原子单位时间平均发射
$1/\tau$个共振光子，
光子能量$\varepsilon=hc/\lambda$。
所以总再辐射功率
$P=(N/\tau)(hc/\lambda)$，得到
\[
 \boxed{N=\frac{P\tau\lambda}{hc}
 \simeq6.7\times10^{9}\approx7\times10^9}.
\]
\excheck 已假定自发辐射主要经题给共振线退激；
若还有非辐射通道，则需要额外分支比。

\atomictermsproblem{129}{247}
锂原子浓度$n=3.6\times10^{16}\,\mathrm{cm^{-3}}$，
温度$T=1500\,\mathrm K$；
其$2P\to2S$共振线$\lambda=671\,\mathrm{nm}$
单位气体体积辐射功率为
$\mathcal P=0.30\,\mathrm{W/cm^3}$。
求共振激发态平均寿命。
\exsolution
热平衡激发态浓度
$n_{2P}=n(g_P/g_S)e^{-hc/(\lambda k_BT)}$，
其中$g_P/g_S=3$。
每个激发原子自发退激功率
$\varepsilon/\tau=hc/(\lambda\tau)$，
故
\[
 \mathcal P=
 \frac{3n\,e^{-hc/(\lambda k_BT)}}{\tau}
 \frac{hc}{\lambda}.
\]
解得
\[
 \boxed{\tau=
 \frac{3n\,hc}{\lambda\mathcal P}
 \exp\!\left[-\frac{hc}{\lambda k_BT}\right]
 \simeq6.6\times10^{-8}\,\mathrm s
 \approx65\,\mathrm{ns}}.
\]
\excheck $n$与$\mathcal P$都以$\mathrm{cm^{-3}}$
和$\mathrm{W/cm^3}$给出，体积单位在比值中抵消。

\atomictermsproblem{130}{247}
氢原子与自身辐射处于热平衡。
求：(a)$T=3000\,\mathrm K$时
从$2P$发射的受激与自发跃迁概率比；
(b)二者相等时的温度。
\exsolution
爱因斯坦系数结合普朗克辐射谱给
\[
 \frac{P_{\rm ind}}{P_{\rm sp}}=
 \frac1{\exp(\hbar\omega/k_BT)-1}.
\]
对$2P\to1S$共振跃迁，
$\hbar\omega=\Delta E=10.2\,\mathrm{eV}$。
(a) 代$T=3000\,\mathrm K$得
\[
 \boxed{\frac{P_{\rm ind}}{P_{\rm sp}}
 \simeq7.3\times10^{-18}}.
\]
(b) 比值等于1要求
$e^{\Delta E/(k_BT)}=2$，
\[
 \boxed{T=\frac{\Delta E}{k_B\ln2}
 \simeq1.71\times10^5\,\mathrm K}.
\]
\excheck 用户底本书末(a)印约$10^{-34}$，
与同一本书所用的$T=3000\,$K及
$10.2\,$eV代入相差约$10^{16}$倍。
本译解保留正确代入值，登记书末数值争议。
(b)的$1.71\times10^5$K是\emph{理想二能级
热平衡辐射模型的形式解}，不能据此断言
实际氢气仍由中性氢原子组成：
如此高温下电离平衡、电离连续谱以及
碰撞过程都必须计入，二能级近似不再自洽。

\atomictermsproblem{131}{247}
一束角频率等于气体原子两能级共振频率$\omega$的光，
通过温度$T$的稀薄原子气体。
考虑受激发射证明吸收系数
$\chi=\chi_0(1-e^{-\hbar\omega/k_BT})$，
其中$\hbar\omega\gg k_BT$，
$\chi_0$为低温极限吸收系数。
\exsolution
设低、上能级占据数密度分别$n_1,n_2$，
简并度$g_1,g_2$；
爱因斯坦系数满足$g_1B_{12}=g_2B_{21}$。
净吸收光功率与
$n_1B_{12}-n_2B_{21}$成正比，
因此
\[
 \chi=\frac{\hbar\omega}{c}
 (n_1B_{12}-n_2B_{21})
 =\frac{\hbar\omega}{c}n_1B_{12}
 \left(1-\frac{n_2g_1}{n_1g_2}\right).
\]
热平衡玻尔兹曼分布
$n_2/n_1=(g_2/g_1)e^{-\hbar\omega/k_BT}$，
且$\hbar\omega\gg k_BT$保证$n_1\simeq n$为总浓度，
所以
\[
 \boxed{\chi=\chi_0(1-e^{-\hbar\omega/k_BT})},
 \quad\chi_0=\frac{\hbar\omega nB_{12}}c.
\]
\excheck 该修正因子来自受激发射抵消一部分吸收，
并非改动单个原子的吸收系数定义。

\atomictermsproblem{132}{247}
汞共振线波长$\lambda=253.65\,\mathrm{nm}$，
激发态平均寿命$\tau=0.15\,\mu\mathrm s$。
求温度$T=300\,\mathrm K$时
多普勒展宽与自然线宽之比的数量级。
\exsolution
最概然热速率
$v_{\rm mp}=\sqrt{2k_BT/m_{\rm Hg}}$，
$m_{\rm Hg}\simeq200.59u$。
以$\Delta\omega_D\sim\omega_0v_{\rm mp}/c$
估计多普勒角频率宽度，
自然线宽$\Gamma\sim1/\tau$，
故
\[
 \frac{\Delta\omega_D}{\Gamma}
 \sim\frac{\omega_0v_{\rm mp}\tau}{c}
 =\frac{2\pi v_{\rm mp}\tau}{\lambda}.
\]
若对两侧速度展宽取全宽$2v_{\rm mp}$，
\[
 \boxed{\frac{\Delta\omega_{D,\mathrm{full}}}
 {\Delta\omega_{\mathrm{nat}}}
 \sim\frac{4\pi\tau v_{\rm mp}}{\lambda}
 \simeq1.2\times10^3\approx10^3}.
\]
\excheck 展宽的标准差、半高全宽与粗估全宽
有数量级为1的系数差；
书末亦仅要求估计$\sim10^3$。

\atomictermsproblem{133}{247}
铁（$Z=26$）的$K_\alpha$特征线波长为
$\lambda_{\rm Fe}=193\,\mathrm{pm}$。
利用莫塞莱定律求铜（$Z=29$）
的$K_\alpha$线波长。
\exsolution
对轻中等原子序数的$K_\alpha$线，
$\omega_\alpha\propto(Z-1)^2$，
故$\lambda_\alpha\propto(Z-1)^{-2}$。
于是
\[
 \boxed{\lambda_{\rm Cu}=
 \lambda_{\rm Fe}\left(\frac{26-1}{29-1}\right)^2
 =193\,\mathrm{pm}\frac{625}{784}
 \simeq154\,\mathrm{pm}}.
\]
\excheck 不能直接用$Z^{-2}$取代
题意所用含屏蔽修正的$(Z-1)^{-2}$。

\atomictermsproblem{134}{247}
利用莫塞莱定律求：
(a)铝和钴的$K_\alpha$线波长；
(b)钒原子K、L层电子结合能之差。
\exsolution
类氢屏蔽近似对K线取有效核电荷$Z-1$，
则
\[
 \Delta E_{K\alpha}=
 \mathrm{Ry}(Z-1)^2\left(1-\frac14\right)
 =\frac34\mathrm{Ry}(Z-1)^2,
 \quad \lambda_{K\alpha}=\frac{hc}{\Delta E_{K\alpha}}.
\]
对铝$Z=13$、钴$Z=27$，
\[
 \boxed{\lambda_{\rm Al}\simeq843\,\mathrm{pm}},\qquad
 \boxed{\lambda_{\rm Co}\simeq180\,\mathrm{pm}}.
\]
钒$Z=23$时，K、L结合能差即
$\Delta E_{K\alpha}$，因此
\[
 \boxed{E_K-E_L=
 \frac34(13.606\,\mathrm{eV})(22)^2
 \simeq4.94\,\mathrm{keV}\approx5\,\mathrm{keV}}.
\]
\excheck 该模型忽略了精细结构与精密屏蔽常数，
适用于本题要求的莫塞莱数量级。
